\documentclass{article}

% Punkte-Variable definieren
\newcounter{points}
\newcommand{\addpoints}[1]{\addtocounter{points}{#1}}

% Punkte-Umgebung für Aufgaben
\newenvironment{task}[1]
  {
    \par \noindent \textbf{Aufgabe #1} \par
    \addpoints{#1}
  }
  {\par}

\begin{document}

\section*{Punktetracking Mobile Robotics}

% Aufgabe mit Punktzahl
\begin{task}{17.5}
17,5 von 20\end{task}

\begin{task}{3}
  Diese Aufgabe ist 3 Punkte wert.
\end{task}

% Gesamtpunktzahl anzeigen
\bigskip
\textbf{Gesamtpunktzahl: \thepoints{}}

\end{document}
